% GENERATED FILE. Edit canonical lessons, then run npm run generate:books.
% canonical-source-hash: fnv1a64:87d6c12be39ada3b
% canonical-lessons: TA-C358-kotpatu, TA-C358-tattuvam, TA-C358-cittantam, TA-C358-kolkai, TA-C358-niruvu

\chapter{Theories and Principles}
\label{ch:ta-theories-and-principles}
\hlchaptermodality{358}

\begin{chapteropening}
\textbf{By the end of this chapter:} \emph{I can talk about a theory, a philosophy and a policy.}

Complete the last lesson of chapter 358: I can talk about a theory, a philosophy and a policy.
\end{chapteropening}

\section[\texorpdfstring{\ta{கோட்பாடு} (kōṭpāṭu)}{கோட்பாடு (kōṭpāṭu)}]{\ta{கோட்பாடு} (kōṭpāṭu) — a theory; a doctrine}

\label{lesson:TA-C358-kotpatu}



\begin{quote}
Before the new one: say the Tamil for likewise, then the Tamil for to differ.
\end{quote}

\subsection*{You'll want to know: \ta{கோட்பாடு}}
\textbf{\ta{கோட்பாடு}} — \emph{kōṭpāṭu} — \textquotedblleft{}a theory; a doctrine\textquotedblright{}.

\begin{quote}
\ta{இந்தக்} \ta{கோட்பாடு} \ta{எனக்குப்} \ta{புரியவில்லை}.
\end{quote}

\begin{quote}
\emph{intak kōṭpāṭu eṉakkup puriyavillai.}
\end{quote}

\begin{quote}
I do not understand this theory.
\end{quote}

\ta{கோட்பாட்டை}, the theory as an object: a noun in -\ta{டு} doubles it to -\ta{ட்டு}- before a case ending.

\begin{cousinweb}
\textbf{\ta{கோட்பாடு}} (\emph{kōṭpāṭu}) is native: \emph{kōḷ}, \textquotedblleft{}taking, opinion, tenet\textquotedblright{}, the noun of \emph{koḷ}, \textquotedblleft{}to take, to hold\textquotedblright{} (DEDR 2151), with \emph{pāṭu}, the noun of \emph{paṭu}: \textquotedblleft{}what is held to be so\textquotedblright{}.
\end{cousinweb}

\begin{grammarlens}[title={What you've built}]
One more word for giving reasons and opinions.
\end{grammarlens}

\subsection*{Your turn}
\begin{itemize}
\raggedright
  \item \emph{Say it:} \emph{kōṭpāṭu}
  \item \emph{Say it:} \emph{kōṭpāṭu}, once more
  \item \emph{Say it:} \emph{intak kōṭpāṭu eṉakkup puriyavillai.} — I do not understand this theory.
  \item \emph{Recall:} say the Tamil for likewise, then the Tamil for to differ, then say \emph{kōṭpāṭu} again
\end{itemize}

\begin{tcolorbox}[breakable,colback=teal!4,colframe=teal!35!black,title={Before you move on}]
What does \ta{கோட்பாடு} mean? (\textquotedblleft{}A theory; a doctrine\textquotedblright{}.) Say it once more.
\end{tcolorbox}
\section[\texorpdfstring{\ta{தத்துவம்} (tattuvam)}{தத்துவம் (tattuvam)}]{\ta{தத்துவம்} (tattuvam) — a philosophy; a principle}

\label{lesson:TA-C358-tattuvam}



\begin{quote}
Before the new one: say the Tamil for to differ, then the Tamil for a theory.
\end{quote}

\subsection*{You'll want to know: \ta{தத்துவம்}}
\textbf{\ta{தத்துவம்}} — \emph{tattuvam} — \textquotedblleft{}a philosophy; a principle\textquotedblright{}.

\begin{quote}
\ta{என்} \ta{அண்ணன்} \ta{கல்லூரியில்} \ta{தத்துவம்} \ta{படிக்கிறான்}.
\end{quote}

\begin{quote}
\emph{eṉ aṇṇaṉ kallūriyil tattuvam paṭikkiṟāṉ.}
\end{quote}

\begin{quote}
My elder brother studies philosophy at college.
\end{quote}

\ta{வாழ்க்கைத்} \ta{தத்துவம்}, a philosophy of life. Plural \ta{தத்துவங்கள்}, principles.

\begin{cousinweb}
\textbf{\ta{தத்துவம்}} (\emph{tattuvam}) is borrowed from Sanskrit \emph{tattva}, \textquotedblleft{}the real state of things, truth\textquotedblright{} (CDIAL 5642). Tamil splits the cluster tv with a vowel: tattuvam.
\end{cousinweb}

\begin{grammarlens}[title={What you've built}]
One more word for giving reasons and opinions.
\end{grammarlens}

\subsection*{Your turn}
\begin{itemize}
\raggedright
  \item \emph{Say it:} \emph{tattuvam}
  \item \emph{Say it:} \emph{tattuvam}, once more
  \item \emph{Say it:} \emph{eṉ aṇṇaṉ kallūriyil tattuvam paṭikkiṟāṉ.} — My elder brother studies philosophy at college.
  \item \emph{Recall:} say the Tamil for to differ, then the Tamil for a theory, then say \emph{tattuvam} again
\end{itemize}

\begin{tcolorbox}[breakable,colback=teal!4,colframe=teal!35!black,title={Before you move on}]
What does \ta{தத்துவம்} mean? (\textquotedblleft{}A philosophy; a principle\textquotedblright{}.) Say it once more.
\end{tcolorbox}
\section[\texorpdfstring{\ta{சித்தாந்தம்} (cittāntam)}{சித்தாந்தம் (cittāntam)}]{\ta{சித்தாந்தம்} (cittāntam) — an ideology; a doctrine}

\label{lesson:TA-C358-cittantam}



\begin{quote}
Before the new one: say the Tamil for a theory, then the Tamil for a philosophy.
\end{quote}

\subsection*{You'll want to know: \ta{சித்தாந்தம்}}
\textbf{\ta{சித்தாந்தம்}} — \emph{cittāntam} — \textquotedblleft{}an ideology; a doctrine\textquotedblright{}.

\begin{quote}
\ta{அவருடைய} \ta{சித்தாந்தம்} \ta{எனக்குப்} \ta{பிடிக்கவில்லை}.
\end{quote}

\begin{quote}
\emph{avaruṭaiya cittāntam eṉakkup piṭikkavillai.}
\end{quote}

\begin{quote}
I do not like his ideology.
\end{quote}

A formal word for a whole system of ideas: \ta{அவருடைய} \ta{சித்தாந்தம்}, his ideology.

\begin{cousinweb}
\textbf{\ta{சித்தாந்தம்}} (\emph{cittāntam}) is borrowed from Sanskrit \emph{siddhānta}, \textquotedblleft{}established end, doctrine\textquotedblright{} (CDIAL 13404), from \emph{siddha}, \textquotedblleft{}settled\textquotedblright{}, and \emph{anta}, \textquotedblleft{}end\textquotedblright{}: a settled conclusion.
\end{cousinweb}

\begin{grammarlens}[title={What you've built}]
One more word for giving reasons and opinions.
\end{grammarlens}

\subsection*{Your turn}
\begin{itemize}
\raggedright
  \item \emph{Say it:} \emph{cittāntam}
  \item \emph{Say it:} \emph{cittāntam}, once more
  \item \emph{Say it:} \emph{avaruṭaiya cittāntam eṉakkup piṭikkavillai.} — I do not like his ideology.
  \item \emph{Recall:} say the Tamil for a theory, then the Tamil for a philosophy, then say \emph{cittāntam} again
\end{itemize}

\begin{tcolorbox}[breakable,colback=teal!4,colframe=teal!35!black,title={Before you move on}]
What does \ta{சித்தாந்தம்} mean? (\textquotedblleft{}An ideology; a doctrine\textquotedblright{}.) Say it once more.
\end{tcolorbox}
\section[\texorpdfstring{\ta{கொள்கை} (koḷkai)}{கொள்கை (koḷkai)}]{\ta{கொள்கை} (koḷkai) — a policy; a principle}

\label{lesson:TA-C358-kolkai}



\begin{quote}
Before the new one: say the Tamil for a philosophy, then the Tamil for an ideology.
\end{quote}

\subsection*{You'll want to know: \ta{கொள்கை}}
\textbf{\ta{கொள்கை}} — \emph{koḷkai} — \textquotedblleft{}a policy; a principle\textquotedblright{}.

\begin{quote}
\ta{அரசின்} \ta{புதிய} \ta{கொள்கையை} \ta{மக்கள்} \ta{வரவேற்றார்கள்}.
\end{quote}

\begin{quote}
\emph{araciṉ putiya koḷkaiyai makkaḷ varavēṟṟārkaḷ.}
\end{quote}

\begin{quote}
People welcomed the new government policy.
\end{quote}

Also a personal principle: \ta{அது} \ta{என்} \ta{கொள்கை}, that is my principle.

\begin{cousinweb}
\textbf{\ta{கொள்கை}} (\emph{koḷkai}) is native Dravidian: DEDR 2151 lists Tamil \emph{koḷ}, \textquotedblleft{}to take, receive, regard\textquotedblright{}, and \emph{koḷkai}, \textquotedblleft{}taking, acceptance\textquotedblright{}, the noun that now means \textquotedblleft{}a principle one holds\textquotedblright{}.
\end{cousinweb}

\begin{grammarlens}[title={What you've built}]
One more word for giving reasons and opinions.
\end{grammarlens}

\subsection*{Your turn}
\begin{itemize}
\raggedright
  \item \emph{Say it:} \emph{koḷkai}
  \item \emph{Say it:} \emph{koḷkai}, once more
  \item \emph{Say it:} \emph{araciṉ putiya koḷkaiyai makkaḷ varavēṟṟārkaḷ.} — People welcomed the new government policy.
  \item \emph{Recall:} say the Tamil for a philosophy, then the Tamil for an ideology, then say \emph{koḷkai} again
\end{itemize}

\begin{tcolorbox}[breakable,colback=teal!4,colframe=teal!35!black,title={Before you move on}]
What does \ta{கொள்கை} mean? (\textquotedblleft{}A policy; a principle\textquotedblright{}.) Say it once more.
\end{tcolorbox}
\section[\texorpdfstring{\ta{நிறுவு} (niṟuvu)}{நிறுவு (niṟuvu)}]{\ta{நிறுவு} (niṟuvu) — to establish; to found}

\label{lesson:TA-C358-niruvu}



\begin{quote}
Before the new one: say the Tamil for an ideology, then the Tamil for a policy.
\end{quote}

\subsection*{You'll want to know: \ta{நிறுவு}}
\textbf{\ta{நிறுவு}} — \emph{niṟuvu} — \textquotedblleft{}to establish; to found\textquotedblright{}.

\begin{quote}
\ta{அவர்} \ta{இந்த} \ta{ஊரில்} \ta{ஒரு} \ta{பள்ளியை} \ta{நிறுவினார்}.
\end{quote}

\begin{quote}
\emph{avar inta ūril oru paḷḷiyai niṟuviṉār.}
\end{quote}

\begin{quote}
He founded a school in this town.
\end{quote}

Past (I …): \ta{நிறுவினேன்}, \emph{niṟuviṉēṉ}.

What is founded takes -\ta{ஐ}. A fact can be established too: \ta{உண்மையை} \ta{நிறுவு}.

\begin{cousinweb}
\textbf{\ta{நிறுவு}} (\emph{niṟuvu}) is native Dravidian: DEDR 3675 lists Tamil \emph{niṟuvu}, \textquotedblleft{}to set up, establish\textquotedblright{}, in the family of \emph{nil}, \textquotedblleft{}to stand\textquotedblright{}: to make something stand.
\end{cousinweb}

\begin{grammarlens}[title={What you've built}]
One more word for giving reasons and opinions.
\end{grammarlens}

\subsection*{Your turn}
\begin{itemize}
\raggedright
  \item \emph{Say it:} \emph{niṟuvu}
  \item \emph{Say it:} \emph{niṟuvu}, once more
  \item \emph{Say it:} \emph{avar inta ūril oru paḷḷiyai niṟuviṉār.} — He founded a school in this town.
  \item \emph{Recall:} say the Tamil for an ideology, then the Tamil for a policy, then say \emph{niṟuvu} again
\end{itemize}

\begin{tcolorbox}[breakable,colback=teal!4,colframe=teal!35!black,title={Before you move on}]
What does \ta{நிறுவு} mean? (\textquotedblleft{}To establish; to found\textquotedblright{}.) Say it once more.
\end{tcolorbox}
